\section{Теоретические сведения и описание программного обеспечения}

\subsection{Цель и задачи работы}

Цель работы --- выявить и изучить фонетические закономерности речевого
сигнала: исследовать влияние параметров дискретизации на качество записи,
акустические свойства гласных и согласных фонем русского языка, эффект
коартикуляции.

Основные задачи работы:
\begin{enumerate}
	\item закрепление основных понятий и характеристик речевого сигнала;
	\item приобретение навыков работы с программой для записи и обработки
		речевого сигнала;
	\item практическое исследование фонетических закономерностей речевого
		сигнала.
\end{enumerate}

\subsection{Характеристики речевого сигнала}

Речь образуется в результате прохождения воздушных потоков через речевой
тракт человека. С акустической точки зрения речевой сигнал является
сложным колебательным процессом, состоящим из множества гармонических
составляющих. Каждая гармоника описывается уравнением
\begin{equation}
	S(t) = A \cdot \cos(F \cdot t + \varphi),
\end{equation}
где $A$ --- амплитуда колебаний; $F$ --- частота ($F = 1/T$); $t$ --- время;
$\varphi$ --- фаза.

Разложение речевого сигнала на гармонические составляющие даёт частотный
спектр --- амплитудно-частотную зависимость. На практике динамика спектра
изображается в виде сонограммы (динамической спектрограммы): по оси
абсцисс откладывается время, по оси ординат --- частота, а степень
затемнения соответствует энергии частотных компонент.

Основными видами изображения речевого сигнала являются:
\begin{itemize}
	\item \textbf{спектр} --- амплитудно-частотная зависимость;
	\item \textbf{осциллограмма} --- амплитудно-временная зависимость.
\end{itemize}

В лингвистическом отношении в речевом сигнале отражаются фонетические
(состав фонем) и просодические (ударение, интонация) аспекты речи.

\subsection{Фонетический состав русской речи (задание 1.1)}

Фонема --- минимальная значимая звуковая единица речи. Акустические
свойства фонемы обусловлены артикуляторными особенностями её
образования: местом образования (сужение артикуляторного тракта) и
способом образования (как образуется звук при прохождении воздуха).

В соответствии со способом образования выделяются:
\begin{itemize}
	\item гласные: [у], [о], [а], [э], [ы], [и];
	\item плавные: [й];
	\item боковые (латеральные): [л], [л'];
	\item носовые: [м], [м'], [н'], [н'];
	\item дрожащие: [р], [р'];
	\item щелевые (фрикативные): [с], [с'], [ш], [щ'], [ф], [ф'], [х],
		[х'], [з], [з'], [ж], [в], [в'];
	\item взрывные: [п], [п'], [т'], [т], [к], [к'], [б], [б'], [д],
		[д'], [г'], [г];
	\item аффрикаты: [ц], [ч'].
\end{itemize}

Также фонемы делятся на \textbf{звонкие} и \textbf{глухие}. Полный
перечень русского языка включает 42 фонемы.

Гласные различаются положением тела языка (верхний/нижний ---
передний/задний) и губ (огубленный/неогубленный): например, [у] ---
верхний, задний, огубленный; [а] --- нижний, задний, неогубленный.

Основные акустические признаки классов фонем, используемые в работе:
\begin{itemize}
	\item \textbf{гласные} --- квазипериодический сигнал с чёткой
		гармонической структурой; основной частотой тона $F_0$ и формантами;
	\item \textbf{шипящие} --- шумоподобный (непериодический) сигнал с
		энергией в определённой полосе частот;
	\item \textbf{взрывные} --- кратковременный взрыв, за которым следует
		переход к соседней фонеме;
	\item \textbf{носовые} --- периодический сигнал с ослабленными
		формантами за счёт резонанса носовой полости.
\end{itemize}

\subsection{Используемое программное обеспечение}

Для выполнения работы использовалось следующее программное обеспечение:
\begin{itemize}
	\item \textbf{Audacity 4.0.1} --- бесплатный редактор звука:
		запись речевого сигнала с микрофона, воспроизведение и
		циклическое воспроизведение (Loop Play) выделенных фрагментов
		осциллограммы, выделение участков, масштабирование,
		экспорт в формат WAV (PCM 16 бит);
	\item \textbf{Python 3.14} со следующими библиотеками:
		\begin{itemize}
			\item \texttt{NumPy} --- массивные вычисления, работа с
				отсчётами сигнала;
			\item \texttt{SciPy} --- ресемплинг (\texttt{resample\_poly}),
				поиск пиков автокорреляционной функции;
			\item \texttt{soundfile} --- чтение и запись WAV-файлов
				разных форматов;
			\item \texttt{Matplotlib} --- построение осциллограмм,
				спектрограмм и спектров.
		\end{itemize}
\end{itemize}

Записи выполнялись в формате WAV (Microsoft PCM), частота дискретизации
16000 Гц, моно, разрядность 16 бит. Анализ осциллограмм, спектров и
спектрограмм, сегментация по энергии, оценка основной частоты тона
автокорреляционным методом выполнялись скриптами на языке Python.
Исходные звукозаписи и скрипты сохранены в структуре каталогов
\texttt{audio/lab1/<номер задания>/} и \texttt{scripts/} соответственно.

Задачи 1.1 и 1.2 выполнены: основные характеристики речевого сигнала и
классификация фонем повторены; запись речевых сигналов и регулировка
уровней освоены (см. раздел~\ref{sec:recording}).
